\documentclass[11pt,a4paper,UTF8]{ctexart}
\usepackage{geometry}
\geometry{left=18mm,right=18mm,top=24mm,bottom=24mm,headheight=22pt,footskip=12mm}
\usepackage{amsmath,amssymb,mathtools,bm,esint,extarrows}
\usepackage{xcolor}
\usepackage{tcolorbox}
\tcbuselibrary{skins,breakable}
\usepackage{fancyhdr}
\usepackage{titlesec}
\usepackage{hyperref}
\usepackage{tikz}
\usepackage{booktabs}
\usepackage{tabularx}

% --- Color Palette (Purple & DeepPink Academic Theme) ---
\definecolor{DeepPurple}{HTML}{4C1D95}
\definecolor{TitlePurple}{HTML}{581C87}
\definecolor{SubPurple}{HTML}{6B21A8}
\definecolor{DeepPink}{HTML}{FF1493}
\definecolor{LightPinkBg}{HTML}{FFF5F8}
\definecolor{LightPurpleBg}{HTML}{F8F5FF}
\definecolor{GoldAccent}{HTML}{D97706}
\definecolor{DarkText}{HTML}{1F2937}

\hypersetup{
    colorlinks=true,
    linkcolor=TitlePurple,
    citecolor=DeepPink,
    urlcolor=SubPurple,
    pdfborder={0 0 0}
}

% --- Typography & Spacing ---
\linespread{1.3}
\setlength{\parskip}{0.35em plus 0.1em minus 0.1em}
\setlength{\parindent}{0pt}

% --- Section Titles Styling ---
\titleformat{\section}
  {\Large\bfseries\color{TitlePurple}}{\thesection}{1em}{}
  [\vspace{1mm}{\color{TitlePurple!30}\hrule height 1pt}]
\titleformat{\subsection}
  {\large\bfseries\color{SubPurple}}{\thesubsection}{0.8em}{}
\titleformat{\subsubsection}
  {\normalsize\bfseries\color{DeepPurple}}{\thesubsubsection}{0.6em}{}

% --- tcolorbox Environments ---
% 1. Theorem / Formula / Method / Analysis Box (Deep Purple)
\newtcolorbox{thmbox}[1][]{
    enhanced,
    breakable,
    colback=LightPurpleBg,
    colframe=TitlePurple,
    arc=3pt,
    boxrule=0pt,
    leftrule=3.5pt,
    left=10pt,
    right=10pt,
    top=8pt,
    bottom=8pt,
    title=#1,
    coltitle=white,
    colbacktitle=TitlePurple,
    fonttitle=\bfseries\small,
    attach boxed title to top left={yshift=-2mm, xshift=4mm},
    boxed title style={boxrule=0pt, arc=2pt}
}

% 2. Solution / Formula / Derivation Box (DeepPink)
\newtcolorbox{solbox}[1][]{
    enhanced,
    breakable,
    colback=LightPinkBg,
    colframe=DeepPink,
    arc=3pt,
    boxrule=0pt,
    leftrule=3.5pt,
    left=10pt,
    right=10pt,
    top=8pt,
    bottom=8pt,
    title=#1,
    coltitle=white,
    colbacktitle=DeepPink,
    fonttitle=\bfseries\small,
    attach boxed title to top left={yshift=-2mm, xshift=4mm},
    boxed title style={boxrule=0pt, arc=2pt}
}

% 3. Learning Goals / Summary Box
\newtcolorbox{targetbox}[1][]{
    enhanced,
    breakable,
    colback=LightPurpleBg,
    colframe=DeepPurple,
    arc=3pt,
    boxrule=1pt,
    left=10pt,
    right=10pt,
    top=8pt,
    bottom=8pt,
    title=#1,
    coltitle=white,
    colbacktitle=DeepPurple,
    fonttitle=\bfseries\small,
    attach boxed title to top left={yshift=-2mm, xshift=4mm},
    boxed title style={boxrule=0pt, arc=2pt}
}

\begin{document}

% ------------------- 讲义封面 (Cover Page) -------------------
\begin{titlepage}
\thispagestyle{empty}
\begin{tikzpicture}[remember picture,overlay]
    \fill[DeepPurple] (current page.north west) rectangle ([yshift=-7cm]current page.north east);
    \draw[DeepPink, line width=3.5pt] ([yshift=-7cm]current page.north west) -- ([yshift=-7cm]current page.north east);
    \draw[GoldAccent, line width=1pt] ([yshift=-7.12cm]current page.north west) -- ([yshift=-7.12cm]current page.north east);
    
    \node[anchor=north east, opacity=0.12, text=white] at ([xshift=-1.5cm, yshift=-0.5cm]current page.north east) {
        \fontsize{70}{70}\selectfont\bfseries 2027
    };
\end{tikzpicture}

\vspace*{0.8cm}
\begin{center}
    {\color{white}\fontsize{26}{32}\selectfont\bfseries 2027 考研数学(二) 核心讲义}\\[0.6cm]
    {\color{white}\fontsize{13}{18}\selectfont\bfseries 全国硕士研究生招生考试 · 统考数学(二) 核心精讲全书}\\[1.5cm]
\end{center}

\vspace{3.5cm}
\begin{center}
    {\color{GoldAccent}\large\bfseries $\blacktriangleright$ 基础强化 · 核心题解 · 考点深水区 $\blacktriangleleft$}\\[0.8cm]
    {\color{TitlePurple}\fontsize{24}{28}\selectfont\bfseries 第二章 数列极限}\\[0.6cm]
    {\color{DeepPink}\large\bfseries —— 体系化理论精讲与公式全景 ——}\\[1.5cm]
\end{center}

\vfill

\begin{center}
\begin{tcolorbox}[
    colback=white,
    colframe=DeepPurple,
    width=0.88\textwidth,
    arc=4pt,
    boxrule=1.2pt,
    halign=center,
    drop shadow=gray!30
]
    \vspace{3mm}
    {\bfseries\large 编著团队：EscoffierZhou}\\[2.5mm]
    {\bfseries\color{TitlePurple} 目标院校：中国科学院大学 (UCAS) 计算机专硕 (22408)}\\[2.5mm]
    {\small\color{gray} 备考铁律：手算真题数字 · 错题白纸重做 · 408 客观题为王}\\[2.5mm]
    {\small 编制版本：2027 届首发版 · \today}
    \vspace{2mm}
\end{tcolorbox}
\end{center}
\vspace{0.8cm}
\end{titlepage}

% ------------------- 目录与页面主体配置 -------------------
\pagestyle{fancy}
\fancyhf{}
\fancyhead[L]{\small\color{gray}2027 考研数学(二) · 核心讲义系列}
\fancyhead[R]{\small\color{TitlePurple}\nouppercase{\leftmark}}
\fancyfoot[C]{\small\color{gray}— 第 \thepage\ 页 —}
\fancyfoot[R]{\small\color{gray}UCAS 22408}
\renewcommand{\headrulewidth}{0.5pt}

\pagenumbering{Roman}
\tableofcontents
\newpage
\pagenumbering{arabic}



\section{1.数列的定义和性质}

\begin{thmbox}[1.数列的定义]

  按照某一法则,对应一个确定的实数$x_n$,按$n$从小到大排得到的一个序列:$x_1, x_2, x_3, \dots, x_n$叫作数列,记为$\{x_n\}$($n \in \mathbb{N}_+$)

  Intuition: 数列可视为自变量为正整数的函数,\textbf{数列一定有无穷多项}

\end{thmbox}
\begin{thmbox}[2.数列的性质]

\textbf{[1]} 数列和子列:从$\{a_n\}: a_1, a_2 \dots a_n$中选取无穷多项组成数列$\{a_{n_k}\}: a_{n_1}, a_{n_2} \dots a_{n_k}$

	     Eg:$\{a_{2k}\}: a_2, a_4, \dots, a_{2k}, \dots \quad \{a_{2k-1}\}: a_1, a_3, \dots, a_{2k-1}, \dots$

	  $\to$数列收敛其任何子列收敛

	  $\to$子列发散则原数列发散

	  $\to$两子列收敛于不同值,数列发散(极限唯一性)

\textbf{[2]} 等差数列:$a_n = a_1 +(n-1)d$	$S_n = \frac{n}{2}[2a_1 +(n-1)d] = \frac{n}{2}(a_1 + a_n)$

\textbf{[3]} 等比数列:$a_n = a_1 q^{n-1}$	  $S_n = \frac{a_1(1-q^n)}{1-q}(q \neq 1)$或$= na_1(q=1)$；等比级数:$1+r+r^2 \dots r^n = \frac{1-r^n}{1-r}(r \neq 1)$

\textbf{[4]} 单调数列:若对所有正整数$n$,有$a_{n+1} \ge a_n(a_{n+1} \le a_n)$,则称数列为单调数列

\textbf{[5]} 严格单调:若对所有正整数$n$,有$a_{n+1} > a_n(a_{n+1} < a_n)$,则称数列为严格单调

\textbf{[6]} 有界数列:对所有正整数$n$,存在正实数$M$,有$|a_n| \le M$,则$\{a_n\}$为有界数列

    	判定方法[1]:找$M$使$|a_n| \le M$

    	判定方法[2]:放缩法  

    	判定方法[3]:找最值  

    	判定方法[4]:基本不等式法

\textbf{[6]} 常见数列和:

  a.$\sum_{k=1}^n k = 1+2+3+\dots+n = \frac{n(n+1)}{2}$

  b.$\sum_{k=1}^n k^2 = 1^2+2^2+3^2+\dots+n^2 = \frac{n(n+1)(2n+1)}{6}$

  c.$\sum_{k=1}^n \frac{1}{k(k+1)} = \frac{1}{1 \times 2} + \frac{1}{2 \times 3} + \dots + \frac{1}{n(n+1)} = \frac{n}{n+1}$(裂项)$\to \lim_{n \to \infty} \sum_{k=1}^n \frac{1}{k(k+1)} = 1$

\textbf{[7]} $\{(1+\frac{1}{n})^n\}$:单调递增有上界,且$\lim_{n \to \infty}(1+\frac{1}{n})^n = e$

\end{thmbox}

\section{2.数列极限的定义}
数列极限的定义[1]:$\{x_n\}$为一数列,若存在常数$a$,[对于任意$\varepsilon > 0$,总存在正整数$N$,使得当$n > N$时,$|x_n - a| < \varepsilon$恒成立],则称$a$为数列的极限

数列极限的定义[2]:数列$\{x_n\}$收敛于$a$,记作$\lim_{n \to \infty} x_n = a$或$x_n \to a$；若不存在$a$,则数列发散

数列和函数极限的区别:

\begin{solbox}[分步推导与答题规范]
  共性:都具有唯一性(所以即使无穷小不能用函数等价无穷小$\to$海涅)

  区别:函数\textbf{局部有界}；数列\textbf{全局有界}(任意一项都有$|a_n| \le M$)

      函数\textbf{局部保号(气泡)}；数列\textbf{后面长保号(射线)}(须存在正整数$N, n>N$有$a_n>0$)
\end{solbox}


\section{3.数列极限的无穷小和无穷大量}
($\varepsilon > 0, X > 0, x > X, |f(x)-A|<\varepsilon$函数)

无穷小:$\lim_{n \to \infty} x_n = a$ $\Leftrightarrow$任意$\varepsilon > 0$,存在$N \in \mathbb{N}_+$,当$n > N$时,恒有$|x_n - a| < \varepsilon$,且$a=0$,称$x_n$是$n \to \infty$时的无穷小量

无穷大:$\lim_{n \to \infty} x_n = \infty$ $\Leftrightarrow$任意$X > 0$,存在$N \in \mathbb{N}_+$,当$n > N$时,恒有$|x_n| > X$,称$x_n$为$n \to \infty$时无穷大


\section{4.数列极限(收敛数列) 的性质}
唯一性:给出数列$\{x_n\}$,若$\lim_{n \to \infty} x_n = a$存在,则$a$是唯一的		(活用$(-1)^n, \sin(\frac{n\pi}{2})$)

有界性:若数列$\{x_n\}$极限存在,则数列$\{x_n\}$有界	$\Leftrightarrow$		<font color=red>数列有界不一定有极限($(-1)^n$)</font>

保号性[1]:(全$\to$局部) 设$\lim_{n \to \infty} x_n = A > \text{常数}b$(或$<\text{常数}b$),则存在$N > 0$,当$n > N$时有$x_n > b$(或$<b$)

保号性[2]:(局$\to$全部) 若数列$\{x_n\}$从某项起有$x_n \ge b(\le)$,且$\lim_{n \to \infty} x_n = a$,则$a \ge b(\le)$,其中$b$为任意常数

脱帽/戴帽法:

\begin{solbox}[分步推导与答题规范]
  脱帽法:$\lim_{n \to \infty} x_n > a$	$\Rightarrow x_n > a$(只能推严格大于,如$x_n = \frac{1}{n} \to \lim_{n \to \infty} x_n = 0$,但实际每一项$x_n > 0$)

  戴帽法:$x_n \ge a(x_n > a)$ $\Rightarrow \lim_{n \to \infty} x_n \ge a$ $\to$ $x_n = \frac{1}{n}$每一项都$>0$,但极限为$0(\ge 0)$
\end{solbox}


\section{5.数列极限计算}

\begin{thmbox}[[1]极限四则运算规则:]

设$\lim_{n \to \infty} x_n = a$,$\lim_{n \to \infty} y_n = b$

$\lim_{n \to \infty}(x_n \pm y_n) = \lim_{n \to \infty} x_n \pm \lim_{n \to \infty} y_n$；

$\lim_{n \to \infty}(x_n - y_n) = \lim_{n \to \infty} x_n - \lim_{n \to \infty} y_n$

$\lim_{n \to \infty} x_n y_n = \lim_{n \to \infty} x_n \cdot \lim_{n \to \infty} y_n$

$\lim_{n \to \infty} \frac{x_n}{y_n} = \frac{\lim_{n \to \infty} x_n}{\lim_{n \to \infty} y_n} (\lim_{n \to \infty} y_n \neq 0)$

\end{thmbox}
\begin{thmbox}[[2]海涅定理(归结原则):]
\textit{}<font color=deeppink>数列特点:题干中有显式通项公式,且内部不包含$n!$或$(-1)^n$等离散特有属性</font>\textit{}

前提: 极限存在的条件下,函数极限(用于计算)和数列极限(采样证明)可以相互转化

操作: 将数列极限直接转化为函数极限$\to$常数/$\infty$/不存在(需要反用解释)

     <font color=red>$\Rightarrow$避坑:$\sin(n\pi)$海涅是$0$,但$\lim$不存在,不能说极限不存在</font>

\end{thmbox}
\begin{thmbox}[[3]夹逼准则(横向展开的,项数太多,所以需要首尾压缩)]

数列特点[1]:巨大的求和$\sum$

数列特点[2]:每一项都在动态变化(尤其是分母)

数列特点[3]:分子分母变量不同阶

原文:若$\{x_n\}, \{y_n\}, \{z_n\}$满足:从某项起,存在$n_0 \in \mathbb{N}$,$n > n_0$时,$y_n \le x_n \le z_n \Rightarrow$且$\lim_{n \to \infty} y_n = a = \lim_{n \to \infty} z_n$,则$\{x_n\}$极限存在且为$a$

\textbf{操作[1]:简单的放大缩小}

(无限项)$n \cdot U_{\min} \le u_1 + u_2 + \dots + u_n \le n \cdot U_{\max} \searrow$

					    求组合极限相等

(有限项)$1 \cdot U_{\min} \le u_1 + u_2 + \dots + u_n \le 1 \cdot U_{\max} \nearrow$

\textbf{操作[2]:不等式}

(三角不等式)$|a \pm b| \le |a| + |b|$,$||a| - |b|| \le |a-b| \Rightarrow |a_1 \pm a_2 \dots \pm a_n| \le |a_1| + |a_2| + \dots + |a_n|$

(均值不等式)$\sqrt{ab} \le \frac{a+b}{2} \le \sqrt{\frac{a^2+b^2}{2}} (a,b \ge 0) \Rightarrow \sqrt[3]{abc} \le \frac{a+b+c}{3} \le \sqrt{\frac{a^2+b^2+c^2}{3}}$

(指数关系)$a \ge b \ge 0$时:[$m > 0$时$a^m \ge b^m$]；[$m < 0$时$a^m \le b^m$]

(分式系) 

  [$0 < a < x < b$]		    [$0 < c < y < d$]

  Eg:$n\pi < x <(n+1)\pi$,$2n < f(x) < 2(n+1)$,则$\frac{2n}{(n+1)\pi} < \frac{f(x)}{x} < \frac{2(n+1)}{n\pi}$

  Eg:$n \le x < n+1$,$2n \le f(x) < 2(n+1)$,则$\frac{2n}{n+1} < \frac{f(x)}{x} < \frac{2n+1}{n} \Rightarrow$取极限$\lim_{n \to \infty} \frac{f(x)}{x} = 2$

(放缩不等式)

  [$\sin x < x < \tan x (0 < x < \frac{\pi}{2})$]	[$\sin x < x (x > 0)$]	[$\arctan x \le x \le \arcsin x (0 \le x \le 1)$]

  [$x < \tan x < \frac{4}{\pi}x (0 < x < \frac{\pi}{4})$	  [$\sin x > \frac{2}{\pi}x (0 < x < \frac{\pi}{2})$]  [$e^x \ge x+1 (\forall x)$]

  [$x-1 \ge \ln x (x>0)$]	        [$\frac{x}{1+x} < \ln(1+x) < x (x>0)$]

\textbf{操作[3]:极值定界}

利用闭区间上连续的函数必有最大值和最小值（代入函数值$f(x_0)$)

\end{thmbox}
\begin{thmbox}[[4]压缩映射原理(纵向展开的,只给了递推关系,动点越来越接近靶心)]

数列特点:专治递推数列$x_{n+1} = f(x_n)$

原理[1]:数列$\{x_n\}$,存在常数$k (0 < k < 1)$,使$|x_{n+1} - a| \le k|x_n - a|$,$n=1,2,\dots$,则$\{x_n\}$收敛于$a$

  解释:该工具的原理:下一次的误差会成比例缩小一圈,最后误差被锁定到0

原理[2]:数列$\{x_n\}$,若$x_{n+1} = f(x_n), n=1,2,\dots$,$f(x)$可导,$a$是$f(x)=x$的唯一解,对$\forall x \in \mathbb{R}$,$|f'(x)| \le k < 1$,则$\{x_n\}$收敛于$a$

  解释:该工具的前提:从递推关系式找到压缩比列k

标准步控SOP(极其重要):

  [步骤一:找靶心]极限解方程  	令$\lim x_n = A \Rightarrow A = f(A)$,算出$A$ $\Rightarrow$(强调算出来的是可能的极限)

  [步骤二:误差方程]作差处理方程      新项差$x_{n+1} - A = f(x_n) - f(A)$

  [步骤三:拉格朗日变型成原理一]      拉格朗日中值:$f(x_n) - f(A) = f'(\xi)(x_n - A)$(其中$x_n < \xi < A$)

  				$\Rightarrow$找在定义域上的上界:证$|f'(\xi)| \le L < 1$

  				(如果你能保证函数$f(x)$的导数在任何地方都满足$\vert{}f'(x)\vert{} \le k < 1$,那么$\vert{}f'(\xi)\vert{}$自然也小于等于$k$)

  [步骤四:计算误差大小变化\}	$|x_{n+1} - A| = |f'(\xi)| \cdot |x_n - A|$ $\Rightarrow$ $|x_{n+1} - A| \le L|x_n - A|$

  [步骤五:递推展开]		 如果找到$|f'(x)|$上界小于$1$,则$|x_{n+1} - A| \le L|x_n - A| \le L^2|x_{n-1} - A| \dots \le L^n|x_1 - A|$

  [步骤六:最后整理] 		$0 < L < 1 \to L^n \to 0 \Rightarrow |x_{n+1} - A| \to 0 \Rightarrow x_n \to A$

Eg:已知数列$\{x_n\}$满足$x_1 = 0$，且$x_{n+1} = \sqrt{2 + x_n}\ (n=1,2,\dots)$证明数列$\{x_n\}$收敛

[步骤一:找靶心] 极限解方程:	    令$\lim_{n \to \infty} x_n = A$,$\therefore$靶心方程:$A = \sqrt{2 + A}$

  解得$A = 2$或$A = -1$,(由于$x_1 = 0$且存在根号，显然$x_n \ge 0$)

[步骤二:误差方程] 作差处理方程:	$x_{n+1} - 2 = \sqrt{2 + x_n} - \sqrt{2 + 2}$

  令核心函数$f(x) = \sqrt{2 + x}$，则上式完美转化为$f(x_n) - f(2)$

[步骤三:拉格朗日变型成原理一] 对$f(x)$在$x_n$与$2$之间使用拉格朗日中值定理:

  $x_{n+1} - 2 = f'(\xi)(x_n - 2)$，其中$\xi$介于$x_n$与$2$之间,现在则需要找到一个L

  >   $f(x) = \sqrt{2 + x}$,$\ \therefore f'(x) = \frac{1}{2\sqrt{2 + x}}$
  >
  >   寻找上界$L$: 因为$x_1 = 0$且递推关系使得所有$x_n \ge 0$，所以介于中间的$\xi \ge 0$
  >
  >   因此，导数值$\vert{}f'(\xi)\vert{} = \frac{1}{2\sqrt{2 + \xi}} \le \frac{1}{2\sqrt{2}} < \frac{1}{2}$

[步骤四:计算误差大小变化] 对步骤三的等式两边取绝对值:

  $\vert{}x_{n+1} - 2\vert{} = \vert{}f'(\xi)\vert{} \cdot \vert{}x_n - 2\vert{} \le \frac{1}{2}\vert{}x_n - 2\vert{}$

[步骤五:递推展开]开始无情降维压缩:

  $\vert{}x_{n+1} - 2\vert{} \le \frac{1}{2}\vert{}x_n - 2\vert{} \le (\frac{1}{2})^2\vert{}x_{n-1} - 2\vert{} \le \dots \le (\frac{1}{2})^n\vert{}x_1 - 2\vert{}$

[步骤六:最后整理]

  由于$0 < \frac{1}{2} < 1$，当$n \to \infty$时，$(\frac{1}{2})^n \to 0$

  由夹逼定理可知，$\lim_{n \to \infty} \vert{}x_{n+1} - 2\vert{} = 0$

  所以原数列极限存在，且$\lim_{n \to \infty} x_n = 2$

\end{thmbox}
\begin{thmbox}[[5]单调有界准则[递推数列$x_{n+1} = f(x_n)$全集]]

原文:单调有界数列必有极限,即若数列$\{x_n\}$单调递增(减)且有上(下)界，则$\lim_{n \to \infty} x_n$存在

  即:IF递增:$x_n \le x_{n+1} \le a$； IF递减:$a \le x_{n+1} \le x_n$

题型:单调有界准则和夹逼拉格朗日(用于$x_{n+1} = f(x_n)$,区别在于对$f(x)$的观察)

\textit{}步骤[1]:求导:先对$f(x)$求导得到$f'(x)$(能量波动形式):\textit{} 

  $f'(x)>0$（能量单向变化)	$\to$\textbf{单调有界准则}

  $f'(x)<0$，$f(x)$震荡（能量左右乱跳)$\to$\textbf{压缩映射定理} 

  (注:这里并非求$f(x_0)$或$f'(\xi)$或上界，单纯看$f'(x)$函数是否恒正/恒负)

  补充:

  

\textit{}步骤[2A]($f'(x)>0$)单调有界准则SOP:\textit{}

  步骤[1]证有界性: 解$A=f(A)$，根据$x_1$和$A$判断上界(一般$A$就是)

  >   其它有界数列判定法[1]:找$M$,使$|a_n| \le M$放缩法 
  >
  >   其它有界数列判定法[2]:找最值
  >
  >   其它有界数列判定法[3]:基本不等式

  步骤[2]证单调性:拉格朗日:$x_{n+1} - x_n = f(x_n) - f(x_{n-1}) = f'(\xi)(x_n - x_{n-1})$

  	       因为$f'(x)>0$，可见$x_{n+1}-x_n$和$x_n-x_{n-1}$同号

  		算一下$x_2-x_1 > 0 (<0)$，所以任意$n$，$x_{n+1}-x_n > 0 (<0)$

  >   其它数列单调证明法[1]:作差法/数学归纳/作商法)

  步骤[3]下结论: 由单调有界准则，极限存在，令$\lim x_n = A$，则$A=f(A)$得极限为$A$

\textit{}步骤[2B]($f'(x)<0$)压缩映射定理SOP:\textit{}

[步骤一:找靶心]极限解方程  	令$\lim x_n = A \Rightarrow A = f(A)$,算出$A$ $\Rightarrow$(强调算出来的是可能的极限)

[步骤二:误差方程]作差处理方程      新项差$x_{n+1} - A = f(x_n) - f(A)$

[步骤三:拉格朗日变型成原理一]      拉格朗日中值:$f(x_n) - f(A) = f'(\xi)(x_n - A)$(其中$x_n < \xi < A$)

			$\Rightarrow$找在定义域上的上界:证$|f'(\xi)| \le L < 1$

			(如果你能保证函数$f(x)$的导数在任何地方都满足$\vert{}f'(x)\vert{} \le k < 1$,那么$\vert{}f'(\xi)\vert{}$自然也小于等于$k$)

[步骤四:计算误差大小变化\}	$|x_{n+1} - A| = |f'(\xi)| \cdot |x_n - A|$ $\Rightarrow$ $|x_{n+1} - A| \le L|x_n - A|$

[步骤五:递推展开]		 如果找到$|f'(x)|$上界小于$1$,则$|x_{n+1} - A| \le L|x_n - A| \le L^2|x_{n-1} - A| \dots \le L^n|x_1 - A|$

[步骤六:最后整理] 		$0 < L < 1 \to L^n \to 0 \Rightarrow |x_{n+1} - A| \to 0 \Rightarrow x_n \to A$

\textit{}步骤[2C]($f'(A)=0$)直接通过靶心公式就可以得出结论:\textit{}

原理:(一阶泰勒)$x_{n+1}-A = f(x_n)-f(A) \approx f'(A)(x_n-A)$	($f'(A)=0$,失效)

   (二阶泰勒)$x_{n+1}-A \approx \frac{1}{2}f''(A)(x_n-A)^2$	

操作[1]:已知$\lim_{n \to \infty} x_n = A$且$f(A)=A, f'(A)=0$求极限$\lim_{n \to \infty} \frac{x_{n+1}-A}{(x_n-A)^2}$

操作[2]:由于动点极度逼近靶心，对$f(x_n)$在$A$处进行二阶泰勒展开：

  
\begin{equation*}
x_{n+1} = f(x_n) = f(A) + f'(A)(x_n-A) + \frac{1}{2}f''(A)(x_n-A)^2 + o((x_n-A)^2)
\end{equation*}


操作[3]:降维:因为$f(A)=A$且$f'(A)=0$，带入后前两项彻底抵消，留下主部：

  
\begin{equation*}
x_{n+1} - A \approx \frac{1}{2}f''(A)(x_n-A)^2
\end{equation*}


  
\begin{equation*}
\lim_{n \to \infty} \frac{x_{n+1}-A}{(x_n-A)^2} = \lim_{n \to \infty} \frac{\frac{1}{2}f''(A)(x_n-A)^2}{(x_n-A)^2} = \frac{1}{2}f''(A)
\end{equation*}


\textit{}*

Eg:已知数列$x_{n+1} = \frac{1}{2}(x_n + \frac{2}{x_n})$，且已知该数列收敛于$\sqrt{2}$求极限$\lim_{n \to \infty} \frac{x_{n+1}-\sqrt{2}}{(x_n-\sqrt{2})^2}$

解:令核心函数$f(x) = \frac{1}{2}(x + \frac{2}{x})= \frac{1}{2}x + x^{-1}$,因此已知极限靶心$A = \sqrt{2}$

步骤1:[验证导数是否为0]求一阶导数：$f'(x) = \frac{1}{2}(1 - \frac{2}{x^2})$

  代入靶心$A = \sqrt{2}$：$f'(\sqrt{2}) = \frac{1}{2}(1 - \frac{2}{2}) = 0$（完美触发二阶收敛定理）

步骤2:[算速度:求二阶导数并代入靶心]:$f''(x) = 2x^{-3} = \frac{2}{x^3}$

  代入靶心$A = \sqrt{2}$,$f''(\sqrt{2}) = \frac{2}{(\sqrt{2})^3} = \frac{2}{2\sqrt{2}} = \frac{\sqrt{2}}{2}$

步骤3:[出结论:根据泰勒降维结论]:$f(x_n) = f(\sqrt{2}) + f'(\sqrt{2})(x_n - \sqrt{2}) + \frac{1}{2}f''(\sqrt{2})(x_n - \sqrt{2})^2 + o((x_n - \sqrt{2})^2)$

带入公式$x_{n+1} = \sqrt{2} + \frac{\sqrt{2}}{4}(x_n - \sqrt{2})^2 + o((x_n - \sqrt{2})^2)$

步骤4:代回原式:

$I = \lim_{n \to \infty} \frac{[\sqrt{2} + \frac{\sqrt{2}}{4}(x_n - \sqrt{2})^2 + o((x_n - \sqrt{2})^2)] - \sqrt{2}}{(x_n - \sqrt{2})^2} = \lim_{n \to \infty} \frac{\frac{\sqrt{2}}{4}(x_n - \sqrt{2})^2 + o((x_n - \sqrt{2})^2)}{(x_n - \sqrt{2})^2} = \lim_{n \to \infty} \left( \frac{\sqrt{2}}{4} + \frac{o((x_n - \sqrt{2})^2)}{(x_n - \sqrt{2})^2} \right)= \frac{\sqrt{2}}{4}$

\textit{}步骤[2D]($f(x)$十分难算):包含超越式/分段函数/$x_{n+1}$难型(不可导)\textit{}

核心:强心捏造出单调性和有界性,然后直接设$\lim_{n \to \infty} x_n = A$,解方程$A = f(A)$

情况(1):[不可导]:手动算前几项,然后脱掉尖点

  $x_{n+1} = |x_n - 1| + 2$，在$x=1$处存在尖点(不可导)		$\to$第一步失效

  SOP: 试算$[x_1=0], [x_2=|0-1|+2=3], [x_3=|3-1|+2=4]$  $\to$绝对值号没必要了，得到$x_{n+1} = x_n + 1$

  [后续:单调递增,没有上界]

情况(2):[超越式]:恶心的都是$e^x$和嵌套,所以可以进行放缩

  $x_{n+1} = x_n e^{-x_n} + \ln(1+x_n^2)$，求导会死人		     $\to$第一步失效

  SOP: 超越式看降放缩:$e^x \ge 1+x, e^{-x_n} \ge 1-x_n$	 	$\Rightarrow x_n e^{-x_n} \ge x_n - x_n^2$；

  		  $\ln(1+x) \le x$                	$\Rightarrow \ln(1+x_n^2) \le x_n^2$

  						 	$\Rightarrow$得到$x_{n+1} - x_n \le 0$

  [后续:回到原式$x_{n+1} = x_n e^{-x_n} + \ln(1+x_n^2)$,两边取极限得$A = A e^{-A} + \ln(1+A^2)$解这个方程，得出$A=0$]

情况(3):[不等式]:通过不等式直接得到了下界

  $x_{n+1} = \frac{1}{2}(x_n + \frac{a}{x_n})$(其中$a>0, x_1>0$)，拉格朗日繁琐$\to$开根号逼近

  操作[1]:有界性(均值不等式):$A+B \ge 2\sqrt{AB}$		$\Rightarrow x_{n+1} = \frac{1}{2}(x_n + \frac{a}{x_n}) \ge \frac{1}{2} \cdot 2\sqrt{x_n \cdot \frac{a}{x_n}} = \sqrt{a}$（有界性)

  操作[2]:作差法:$x_{n+1} - x_n = \frac{1}{2}(\frac{a}{x_n} - x_n) = \frac{a - x_n^2}{2x_n}$因为$x_n \ge \sqrt{a} \Rightarrow x_n^2 \ge a \Rightarrow x_{n+1} - x_n \le 0$（单调性)

  操作[3]:单调有界准则:$A = \frac{1}{2}(A + \frac{a}{A}) \Rightarrow A = \sqrt{a}$

  [后续:回到原式$x_{n+1} = \frac{1}{2}(x_n + \frac{a}{x_n})$,两边取极限得$A = \frac{1}{2}(A + \frac{a}{A})$，解得$A = \sqrt{a}$]

\end{thmbox}
\begin{thmbox}[[6](柯西视角)${x_n}$的收敛于$a$的速度]
\textbf{原文:} 设数列$\{x_n\}, \{y_n\}$在$n \to \infty$过程中同时趋于$A$

(记$U_n=|x_n-a|, V_n=|y_n-a|$，$I = \lim_{n \to \infty} \frac{U_n}{V_n}$，[$n \to \infty, U_n, V_n$都是无穷小量])

若$I=0$，说明$x_n$的收敛速度比$y_n$快	

若$I=k\ (k>0, \text{常数})$，说明$x_n$的收敛速度是$y_n$的$k$倍

若$I=\infty$，说明$x_n$的收敛速度比$y_n$慢

(标准数列极限无法体现出收敛速度$(\lim x_n=a \Leftrightarrow \forall \varepsilon>0, \exists N, n>N, |x_n-a|<\varepsilon) \Leftarrow$少了$\frac{|x_n-a|}{|y_n-a|}=V$)

  书上仅对定义做出了要求，数学符号表示极限 ($\varepsilon$中已取$k$或$N$))

补充[1]:无穷多项数列比大小:	$\ln^\alpha n \ll n^\beta \ll a^n \ll n! \ll n^n\ (a>0, \beta>0, a>1)$

补充[2]:无穷大失效: 	      倒带换元，$t=\frac{1}{n}$，泰勒

<font color=red>补充[3]:$x_{n+1}=f(x_n)$，$y_{n+1}=g(y_n)$，都收敛于$A$,比较$\frac{U_n}{V_n} = \frac{x_n-A}{y_n-A} \Rightarrow$比$f'(A)$和$g'(A)$谁小谁快</font>

\end{thmbox}
\begin{thmbox}[[7](黎曼积分)定积分与数列极限$\Rightarrow$(夹逼拓展) 黎曼积分 (巨大的求和符号)]
\textbf{夹逼拓展:} 

情况[1]$n$的最高次幂和游走变量$i$\textbf{不同阶}做法: 用夹逼:不提取,用Max/Min找上下界（夹逼)

情况[2]$n$的最高次幂和游走变量$i$\textbf{同阶}做法:   用积分:提取$\frac{1}{n}$，剩下必然是$\frac{i}{n}$

步骤[1]:判断齐次性:齐次则启动黎曼积分

步骤[2]:强制析出绝对纯净的微元$\frac{1}{n}$(确立$dx$）

  无论原式多恶心,强行在整个$\sum$内部提取出一个$\frac{1}{n}$

  提完$\frac{1}{n}$后,剩下的表达式里,绝对不允许存在任何落单的$n$,所有的$n$必须且只能作为分母去和$i$强行绑定

步骤[3]:整体组装游走方$\frac{i}{n}$(确立$f(x)$) 

  将剩余表达式中所有的结构强制化为$\frac{i}{n}$的形式,然后将$\frac{i}{n}$整体无脑替换为变量$x$,此时剩余的表达式即为连续函数$f(x)$

步骤[4]:极限映射积分上下界(确立$\int_a^b$)

下界$a$:将求和的第一项指标$i_{\text{start}}$代入，取极限$a = \lim_{n \to \infty} \frac{i_{\text{start}}}{n}$(若从$1$开始，通常$a=0$）

上界$b$:将求和的最后一项指标$i_{\text{end}}$代入，取极限$b = \lim_{n \to \infty} \frac{i_{\text{end}}}{n}$(若加到$n$，通常$b=1$；若加到$2n$，则$b=2$）

步骤[5]:计算[原数列极限直接等价转化为定积分：$I = \int_a^b f(x) dx$]

补充[1]:项数不随$n$增长:$\lim_{n \to \infty} \sum_{i=1}^{10} \frac{1}{n}f(\frac{i}{n})$

  这里总共只有$10$项（常数），提取$\frac{1}{n}$后，极限直接为$0$必须有无穷多项才能积出面积

补充[2]:项数暴增超载:$\sum_{i=1}^{n^2}$

  项数增长是$n^2$级别，游走方$\frac{i}{n}$会冲向无穷大，黎曼积分上下界撕裂失效，通常结果发散或为$0$

Eg:$\lim_{n \to \infty} \sum_{i=1}^n \frac{1}{n}f(\frac{i}{n})$ $\to$ $a = \lim_{n \to \infty} \frac{1}{n} = 0$,$b = \lim_{n \to \infty} \frac{n}{n} = 1$

  特例:$\lim_{n \to \infty} \sum_{i=1}^{10} \frac{1}{n}f(\frac{i}{n})$没有$i$，$i$就是常量，直接$\frac{1}{n} \cdot f(\frac{10}{n}) = 0$，不用积分

  特例:$\sum_{i=1}^{n^2}$因此循环次数必须为$n$同阶的式子，否则都是 0（黎曼积分失效)

\textit{}Eg:$\lim_{n \to \infty} (\frac{n}{n^2+1^2} + \frac{n}{n^2+2^2} + \dots + \frac{n}{n^2+n^2})$\textit{}

\textbf{解:} 游走变量$1^2, 2^2 \dots n^2$，所以$I = \frac{n}{n^2+i^2} \Rightarrow \sum_{i=1}^n \frac{n}{n^2+i^2}$($i$同阶,$n$同阶)

$I = \frac{n}{n^2+i^2} = \frac{1/n}{1+(i/n)^2} \cdot \frac{1}{n} \Rightarrow \lim_{n \to \infty} \frac{1}{n} \sum_{i=1}^n \frac{1}{1+(i/n)^2} = \int_0^1 \frac{1}{1+x^2} dx = \arctan(1) - \arctan(0) = \frac{\pi}{4}$

\textit{}Eg:$\lim_{n \to \infty} (\frac{1}{n^2+n+1} + \frac{2}{n^2+n+2} + \dots + \frac{n}{n^2+n+n})$\textit{}

\textbf{解:}$S_n = \sum_{i=1}^n \frac{i}{n^2+n+i} \Rightarrow$提$\frac{1}{n}$:$S_n = \frac{1}{n} \sum_{i=1}^n \frac{n \cdot i}{n^2+n+i} = \frac{1}{n} \sum_{i=1}^n \frac{i/n}{1+1/n+i/n^2}$($i$的幂次$\neq n$的最高次，积分失效$\to$夹逼)

下界:分母最大即$S_n = \sum_{i=1}^n \frac{i}{n^2+2n} = \frac{1}{n^2+2n} \sum_{i=1}^n i = \frac{1}{n^2+2n} \frac{n(n+1)}{2}$(求和公式)$= \frac{n^2+n}{2n^2+4n}$，$\lim = \frac{1}{2}$

上界:分母最小即$S_n = \sum_{i=1}^n \frac{i}{n^2+n+1} = \frac{1}{n^2+n+1} \sum_{i=1}^n i = \frac{1}{n^2+n+1} \frac{n(n+1)}{2}$(求和公式)$= \frac{n^2+n}{2n^2+2n+2}$，$\lim = \frac{1}{2}$
$\Rightarrow$相同，证毕

\end{thmbox}
\begin{thmbox}[[8]Oth(Stolz定理)处理纯离散未定式]

\textit{}<font color=deeppink>数列特点:一定是分式结构$\frac{x_n}{y_n}$,[分子通常是无法求导的离散累加和(如$\sum$)],[分母单调趋于$\infty$]</font>\textit{}

\textit{}定理内容($\frac{*}{\infty}$型):\textit{} 设$\{y_n\}$严格单调递增,且$\lim_{n \to \infty} y_n = +\infty$,若$\lim_{n \to \infty} \frac{x_{n+1}-x_n}{y_{n+1}-y_n} = A$(A可为实数、$\pm\infty$),则$\lim_{n \to \infty} \frac{x_n}{y_n} = A$

\textbf{SOP(我才不用斯托尔兹降维法):直接夹逼}

Eg:求极限$\lim_{n \to \infty} \frac{1 + \frac{1}{2} + \dots + \frac{1}{n}}{\ln n}$

找下界(积分卡小):$\sum_{i=1}^n \frac{1}{i} > \int_1^{n+1} \frac{1}{x} dx = \ln(n+1)$

找上界(积分卡大):$\sum_{i=1}^n \frac{1}{i} < 1 + \int_1^n \frac{1}{x} dx = 1 + \ln n$


\begin{equation*}
\frac{\ln(n+1)}{\ln n}=1 < \frac{1 + \frac{1}{2} + \dots + \frac{1}{n}}{\ln n} < \frac{1 + \ln n}{\ln n}=1
\end{equation*}

\end{thmbox}

\end{document}
